\section{Related Work}


Petersen et al. introduced Differentiable Logic Gate Networks (DLGNs),
enabling gradient-based optimization of Boolean gate functions while
keeping randomly initialized connections fixed during training
\cite{petersen2022}. In~\cite{petersen2024}, the authors extended DLGNs
to convolutional architectures, exploiting image structure through shared
logic trees, spatial sampling, pooling, and residual initialization. Other
works extend DLGNs to recurrent models \cite{buehrer2025recurrent},
improve output aggregation for larger classification problems
\cite{brandle2025}, or develop interpretable and truth-table-based logic
architectures \cite{yue2024,benamira2022,benamira2024}. Related
differentiable logic models have also been explored for cellular automata
\cite{mordvintsev2020,miotti2025} and circuit representation learning
\cite{li2022deepgate}. However, these works primarily focus on extending
the architecture and applicability of logic-based networks, without exploring the connectivity problem.
% rather than isolating how a deliberately structured \emph{fixed} predecessor graph should be constructed under the same architecture and gate budget.


Regarding DLGN connectivity, Mommen et al. optimize distributions over
candidate predecessors, allowing gate inputs to change during training
\cite{mommen2025}. Similarly, LILogic Net employs sparse Top-$K$ routing
to learn compact network connectivity \cite{fojcik2026}, while BitLogic
studies connectivity together with encoding, fan-in, node
parameterization, and output-head design \cite{buehrer2026}. These works
demonstrate that selecting better predecessors can significantly improve
logic networks. However, learned routing requires processing candidate
connections and maintaining additional routing parameters and
intermediate activations during training, increasing both execution time
and memory requirements, even though the hardened network ultimately
retains only a small number of inputs per gate. In contrast, our approach
constructs the connectivity \emph{offline} and keeps it fixed during gate
training, avoiding an additional routing optimization problem while
preserving the declared depth, width, and gate budget.


Additionally, several works focus on improving the optimization of the
gate functions themselves. Mind the Gap reduces the discrepancy between
continuous training and discrete inference \cite{yousefi2025}. Light
DLGNs redesign the gate parameterization to improve training efficiency
\cite{ruttgers2025}, while WARP introduces a Walsh-based Boolean
representation \cite{gerlach2026}. The works in
\cite{kim2026,damera2026} further explore polynomial formulations to
analyze or replace conventional soft gate mixtures. These methods
optimize \emph{what function} a gate computes from its assigned inputs,
whereas our work focuses on \emph{which signals} are made available to
that gate.